%!TEX root = ../main.tex
\section{Overview of the problem/evaluation}


\subsection{Network traffic analysis}
The network traffic analysis is an important part for developing control 
schemes and distinguish a malicious packet from a normal one. 
Prediction in network traffic is important and there are two categories: 
long and short period predictions. The long period prediction allows more 
detailed forecast, it helps to determine future capacity needs and allows 
better planning and decision taking. Short period predictions can be used 
for congestion control and a better resource management. \cite{network}
\\Nowadays network traffic analysis has become very important 
for monitoring the traffic. Preaviously, only a few newtworks devices were kept 
monitored since the bandwidth was just about 100 Mbps. Now the speed
has increased to more than 1 Gbps and more networks have been created, 
so it is clear how this science has to keep evolving.
There are various techniques for network traffic analysis, the typical approach 
consists of a first processing and inspection or the traffic and then an analysis. 
\subparagraph{Network Traffic Processing} One of the techniques that allow processing a 
large amount of data is anomaly detection, this technique aims to identify suspicious or unusual
activity or behavior. There is a model that rapresents the 'normal network activity' and then all 
network flow is analised based on the model, it monitors traffic load in order to recognise the 
triggers to a significant traffic volume change. These systems are built in two phases: the training 
phase (on a constructed model) and the testing phase (on current traffic). Most of these systems 
detect anomalies using machine or deep learning techniques.
\\Another technique is the signature-based network inspection which mostly uses DPI (Deep Packet Inspection) techniques. 
This tecniques is now used in different layer of the OSI (Open Systems Interconnection), 
unlike earlier which was only used on packet headers. Nowadays protocols are so complex and obfuscated 
that it is necessary to inspect all layers. Using DPI allows to identify different styles of 
content and streaming them based on their type; the same suspicious and malicious traffic is identified and then analised.
However, these DPI based approaches are only able to process unencrypted data because they extract important information from the 
network packet payload content. These tecniques expect unencrypted data which nowadays is not realistic so other 
methods have been investigated. \cite{ntaApplications}
\subparagraph{Network Encryption Analysis}

\subsection{Information Leakage Vectors}

\subsection{The AI/ML Threat Multiplier}

